\documentclass[10pt,a4paper]{article}
\usepackage[margin=20mm,headheight=15pt,headsep=7mm,footskip=10mm]{geometry}
\usepackage{fontspec,amsmath,amssymb,mathtools,microtype}
\setmainfont{lmroman10-regular.otf}[BoldFont=lmroman10-bold.otf,ItalicFont=lmroman10-italic.otf,BoldItalicFont=lmroman10-bolditalic.otf]
\setsansfont{lmsans10-regular.otf}[BoldFont=lmsans10-bold.otf,ItalicFont=lmsans10-oblique.otf,BoldItalicFont=lmsans10-boldoblique.otf]
\setmonofont[Scale=.86]{lmmono10-regular.otf}[ItalicFont=lmmono10-italic.otf]
\usepackage{ragged2e}\usepackage{xcolor,booktabs,longtable,array,graphicx,enumitem,fancyhdr,needspace,titlesec,hyperref,xurl}
\definecolor{AtlasNavy}{HTML}{18334C}\definecolor{AtlasTeal}{HTML}{116D73}\definecolor{AtlasGray}{HTML}{52616D}
\hypersetup{colorlinks=true,linkcolor=AtlasTeal,urlcolor=AtlasTeal,pdfcreator={Complexity Atlas},unicode=true}
\urlstyle{same}\setlength{\parindent}{0pt}\setlength{\parskip}{4.8pt}\setlength{\emergencystretch}{3em}
\renewcommand{\baselinestretch}{1.08}\setcounter{secnumdepth}{0}\setlist{nosep,leftmargin=1.5em}
\titleformat{\section}{\Large\bfseries\sffamily\color{AtlasNavy}\raggedright}{}{0em}{}
\titleformat{\subsection}{\large\bfseries\sffamily\color{AtlasTeal}\raggedright}{}{0em}{}
\titleformat{\subsubsection}{\normalsize\bfseries\color{AtlasNavy}\raggedright}{}{0em}{}
\titlespacing*{\section}{0pt}{14pt}{6pt}\titlespacing*{\subsection}{0pt}{10pt}{5pt}
\pagestyle{fancy}\fancyhf{}\fancyhead[L]{\small\sffamily\color{AtlasGray}COMPLEXITY ATLAS}
\fancyhead[R]{\small\sffamily\color{AtlasGray}ENGLISH EDITION}\fancyfoot[C]{\small\thepage}\renewcommand{\headrulewidth}{.3pt}
\newcommand{\AtlasMath}[1]{\par\begingroup\centering\setbox0=\hbox{$\displaystyle #1$}\ifdim\wd0>1.4\linewidth\typeout{ATLAS WIDE FORMULA: \the\wd0}\fi\ifdim\wd0>\linewidth\resizebox{\linewidth}{!}{\box0}\else\box0\fi\par\endgroup}
\newcommand{\AtlasInlineMath}[1]{\begingroup\setbox0=\hbox{$#1$}\ifdim\wd0>.4\linewidth\endgroup\AtlasMath{#1}\else\endgroup$#1$\fi}
\pdfstringdefDisableCommands{\def\AtlasInlineMath#1{#1}}
\DeclareMathOperator{\polylog}{polylog}\DeclareMathOperator{\dist}{dist}\DeclareMathOperator{\diam}{diam}
\DeclareMathOperator{\rank}{rank}\DeclareMathOperator{\tr}{tr}\DeclareMathOperator{\diag}{diag}
\DeclareMathOperator{\prox}{prox}\DeclareMathOperator{\KL}{KL}\DeclareMathOperator{\Var}{Var}
\providecommand{\eps}{\varepsilon}\providecommand{\R}{\mathbb R}\providecommand{\E}{\mathbb E}
\providecommand{\mathbbm}[1]{\mathbf{#1}}\providecommand{\mathds}[1]{\mathbf{#1}}
\providecommand{\llbracket}{[\![}\providecommand{\rrbracket}{]\!]}
\begin{document}
\raggedright

{\LARGE\bfseries\sffamily\color{AtlasNavy} Finite-sum bilevel optimization: a lower bound under linear-iteration restrictions.\par}\vspace{5pt}
\textbf{Setting ID.} \texttt{BL\_\allowbreak{}SPEC\_\allowbreak{}finite\_\allowbreak{}derivative\_\allowbreak{}linear}

Bilevel optimization

Solve the bilevel problem \AtlasInlineMath{\min_x\Phi(x)}, where \AtlasInlineMath{\Phi(x)=f(x,y^*(x))} and \AtlasInlineMath{y^*(x)=\operatorname*{argmin}_y g(x,y)} is the unique global lower-level solution. The upper- and lower-level variables are \AtlasInlineMath{x\in\mathbb{R}^{d_x}} and \AtlasInlineMath{y\in\mathbb{R}^{d_y}}. Seek approximate stationarity of the reduced objective under the stated error criterion and guarantee. Assumptions on \AtlasInlineMath{f}, \AtlasInlineMath{g}, and \AtlasInlineMath{\Phi} are specified separately. Here \AtlasInlineMath{f=n_f^{-1}\sum_if_i} and \AtlasInlineMath{g=n_g^{-1}\sum_jg_j}.

\textbf{Status.} Bounds missing

\textbf{TASK tags.} task family: Bilevel

\textbf{MODEL tags.} convexity: Lower objective strongly convex in its lower variable; structure: Finite sum

\textbf{ORACLE tags.} interface family: Derivative vector

\textbf{ALGORITHMS tags.} restriction: Linear span

\textbf{OUTPUT tags.} object: Outer point; support: Restricted iteration sequence

\textbf{METRIC tags.} criterion: Reduced-objective gradient norm

\textbf{GUARANTEE tags.} aggregation: Per instance; mode: Expectation

\textbf{RESOURCE tags.} budget: Fixed worst-case budget

\needspace{7\baselineskip}\section{1. Setting specification}
\needspace{4\baselineskip}\subsection{Task}
Solve the bilevel problem \AtlasInlineMath{\min_x\Phi(x)}, where \AtlasInlineMath{\Phi(x)=f(x,y^*(x))} and \AtlasInlineMath{y^*(x)=\operatorname*{argmin}_y g(x,y)} is the unique global lower-level solution. The upper- and lower-level variables are \AtlasInlineMath{x\in\mathbb{R}^{d_x}} and \AtlasInlineMath{y\in\mathbb{R}^{d_y}}. Seek approximate stationarity of the reduced objective under the stated error criterion and guarantee. Assumptions on \AtlasInlineMath{f}, \AtlasInlineMath{g}, and \AtlasInlineMath{\Phi} are specified separately. Here \AtlasInlineMath{f=n_f^{-1}\sum_if_i} and \AtlasInlineMath{g=n_g^{-1}\sum_jg_j}.

\needspace{4\baselineskip}\subsection{Admissible model}
Each \AtlasInlineMath{f_i} has \AtlasInlineMath{L_F}-Lipschitz joint gradient. Each \AtlasInlineMath{g_j} is twice continuously differentiable and \AtlasInlineMath{\mu_y}-strongly convex in \AtlasInlineMath{y}. This lower-bound class imposes no additional uniform Lipschitz budgets on function values, Hessians, or third derivatives.

The points \AtlasInlineMath{x_0,y_0,v_0} are public and \AtlasInlineMath{\Phi(x_0)-\inf\Phi\le\Delta}. No additional bound is imposed on the initial inner or adjoint distance.

\needspace{4\baselineskip}\subsection{Oracle interface}
A query selects \AtlasInlineMath{i,j,x,y,v} and returns the exact joint response \AtlasInlineMath{(\nabla f_i,\nabla g_j,g_{j,yy}v,g_{j,xy}v)}, with every component evaluated at \AtlasInlineMath{(x,y)}. Components may be queried repeatedly.

\needspace{4\baselineskip}\subsection{Output}
Return an outer point from the restricted iteration sequence. An arbitrary terminal output outside the permitted span cannot enlarge the algorithm class.

\needspace{4\baselineskip}\subsection{Loss or accuracy criterion}
\AtlasInlineMath{\|\nabla\Phi(\widehat{x})\|_2^2}.

\needspace{4\baselineskip}\subsection{Resource accounting}
The quantity \AtlasInlineMath{N_{\mathrm{tuple}}} counts joint responses, each costing one unit under a fixed worst-case budget. Four primitive calls can simulate one joint response, but this simulation does not remove the algorithm’s linear-iteration restriction.

\needspace{4\baselineskip}\subsection{Algorithm views and success guarantees}
Restrict to the linear bilevel iteration class in Dagréou et al., Definition 4.2. Component indices may be randomized, but the following history-span restrictions must hold:

\AtlasInlineMath{y_{t+1}\in y_0+\operatorname{span}\{\nabla_yg_{j_s}(x_s,y_s):s\le t\}}; 
\AtlasMath{v_{t+1}\in v_0+\operatorname{span}\{g_{j_s,yy}(x_s,y_s)v_s+\nabla_yf_{i_s}(x_s,y_s):s\le t\}}
; 
\AtlasMath{x_{t+1}\in x_0+\operatorname{span}\{g_{j_s,xy}(x_s,y_s)v_s+\nabla_xf_{i_s}(x_s,y_s):s\le t\}}
.

For every instance, \AtlasInlineMath{\mathbb{E}\|\nabla\Phi(\widehat{x})\|_2^2\le\varepsilon^2}.

\needspace{7\baselineskip}\section{2. Notation and parameters}
\AtlasInlineMath{x,y}: The upper- and lower-level variables. \textbf{Domain:} \AtlasInlineMath{x\in\mathbb R^{d_x},y\in\mathbb R^{d_y}}.

\AtlasInlineMath{f,g}: The original upper- and lower-level real-valued functions. \textbf{Domain:} \AtlasInlineMath{f,g:\mathbb R^{d_x}\times\mathbb R^{d_y}\to\mathbb R}.

\AtlasInlineMath{\widehat x,x_0,y_0}: Output outer point, public outer initial point, and public lower-level initial point \textbf{Domain:} Each lies in its corresponding variable space.

\AtlasInlineMath{\Phi}: The reduced objective. \textbf{Domain:} Defined in the task specification.

\AtlasInlineMath{y^*}: Mapping to the unique lower-level global minimizer \textbf{Domain:} Constraints specified under TASK.

\AtlasInlineMath{\nabla,\|\cdot\|}: Euclidean gradient and vector norm; matrix and higher-derivative norms are the induced Euclidean operator norms

\AtlasInlineMath{\mathbb E,\Pr,\inf,\sup}: Expectation, probability, infimum, and supremum, with aggregation order and randomness specified by the guarantee.

\AtlasInlineMath{\mathcal O,\Omega,\Theta,\widetilde{\mathcal O}}: Asymptotic upper, lower, matching-order, and declared-logarithm upper bounds; fixed parameters are specified in each evidence record

\AtlasInlineMath{c,C}: Positive absolute constants or functions of fixed budgets, with dependencies specified by the corresponding result \textbf{Domain:} \AtlasInlineMath{0<c,C<\infty}.

\AtlasInlineMath{\varepsilon}: Target accuracy. Whether error is squared is determined by the metric and guarantee formulas. \textbf{Domain:} \AtlasInlineMath{\varepsilon>0}.

\AtlasInlineMath{d_x}: Upper-level dimension \textbf{Domain:} Positive integer.

\AtlasInlineMath{d_y}: Lower-level dimension \textbf{Domain:} Positive integer.

\AtlasInlineMath{L_F}: A joint Lipschitz bound on each upper-level component’s full gradient. \textbf{Domain:} \AtlasInlineMath{L_F>0}.

\AtlasInlineMath{\mu_y}: The uniform strong-convexity modulus in the lower-level variable. \textbf{Domain:} \AtlasInlineMath{\mu_y>0}.

\AtlasInlineMath{\Delta}: Public upper bound on the initial reduced-objective gap \textbf{Domain:} \AtlasInlineMath{\Delta\ge0}.

\AtlasInlineMath{\Phi(x_0)-\inf_x\Phi(x)\le\Delta}

\AtlasInlineMath{n_f}: Number of upper-level finite-sum components \textbf{Domain:} Positive integer.

\AtlasInlineMath{n_g}: Number of lower-level finite-sum components \textbf{Domain:} Positive integer.

\AtlasInlineMath{g_{xy},g_{yy}}: Mixed lower-level Hessian block and lower-variable Hessian block \textbf{Domain:} \AtlasInlineMath{g_{xy}\in\mathbb R^{d_x\times d_y},g_{yy}\in\mathbb R^{d_y\times d_y}}.

\AtlasInlineMath{g_{xy}=\mathrm D_y\nabla_xg,\quad g_{yy}=\mathrm D_y\nabla_yg}

\AtlasInlineMath{f_i,g_j}: Upper- and lower-level component functions in the finite sums \textbf{Domain:} \AtlasInlineMath{1\le i\le n_f,\ 1\le j\le n_g}.

\AtlasInlineMath{i,j}: Upper- and lower-level component indices queried by the algorithm

\AtlasInlineMath{v_0,v_t}: Public initial adjoint vector and the algorithm's adjoint iterates. \textbf{Domain:} \AtlasInlineMath{v_0,v_t\in\mathbb R^{d_y}}.

\AtlasInlineMath{g_{j,xy},g_{j,yy}}: Cross-Hessian and lower-variable Hessian blocks of lower component \AtlasInlineMath{j}.

\AtlasInlineMath{t,s,\operatorname{span}}: These denote the current iteration, a historical iteration, and a linear span. Enforce the three algorithm-restriction formulas separately.

\AtlasInlineMath{N,Q}: This is the total call budget under the stated resource projection. If only one symbol is used, the two names are aliases for that same budget. \textbf{Domain:} RESOURCE specifies which operations are summed; no additional unit cost is introduced.

\AtlasInlineMath{N_{\rm tuple}}: Calls returning the joint tuple of two-level gradients and derivative-vector products. \textbf{Domain:} Nonnegative integer.

\needspace{12\baselineskip}\section{3. Complexity bracket}
\needspace{4\baselineskip}\subsection{Compared quantities}
\needspace{4\baselineskip}\subsubsection{Q1 - source scoped metric resource statement}
\textbf{Metric contract.} E||grad Phi(x\_hat)||\textasciicircum{}2 <= epsilon\textasciicircum{}2.

\textbf{Resource contract.} metric: Q\_tuple, the number of calls to the joint tuple in source Eq.(11).; budget: Fixed worst-case call budget.; mapping: One tuple can be simulated by four separate declared component operations, but arbitrary separate-operation algorithms need not satisfy Definition 4.2.

This is the native source metric/resource guarantee. Joint precision, block size, runtime and oracle components stay separate; do not infer a scalar optimal query complexity.

\needspace{4\baselineskip}\subsection{R1 - BM\_BOUND\_finite\_linear\_native\_lower source window}
\textbf{Region.} \AtlasInlineMath{\varepsilon^2\le\Delta L_F/(1000n_f)}. Retain the source smoothness class and the linear-iteration restriction in Definition 4.2. The hard-instance dimension may grow as specified in Theorem 2.

\textbf{Quantity.} Q1

\textbf{Resource.} Joint resource budget

\textbf{Status.} Upper bound missing

\textbf{Lower bound.} B1

\textbf{Upper bound.} No admitted bound for this comparison

\needspace{7\baselineskip}\section{4. Bounds and supporting evidence}
The following conditions apply to every bound in this profile, in addition to each bound's stated region and algorithm scope.

\textbf{Dimension conditions.} The lower-bound hard-instance dimension may grow with \AtlasInlineMath{\Delta,L_F,n_f,1/\varepsilon}. Here \AtlasInlineMath{L_F} is the smoothness budget of the upper-level component gradients. The embedded reduced objective inherits that scale; this does not add a reduced-smoothness assumption to the full model class.

\textbf{Resource mapping.} Charge the joint tuple calls of source equation (11), retaining that restricted model.

\needspace{4\baselineskip}\subsection{B1 - Finite-sum zero-chain embedded through g(y,x)=||y-x||\textasciicircum{}2/2}
\textbf{Admission.} Included in the bracket

\textbf{Direction.} lower

\textbf{Witness or lower-bound scope.} The stated restricted randomized class

\AtlasMath{\Omega(\sqrt{n_f}\,\varepsilon^{-2})}
\textbf{Sources.} \href{https://proceedings.mlr.press/v238/dagreou24a/dagreou24a.pdf}{[1]}

\textbf{Constants.} Hold the stated regularity and initialization budgets fixed. Joint optimality is not claimed for parameters whose dependence is untracked.

\textbf{Source location.} \href{https://proceedings.mlr.press/v238/dagreou24a/dagreou24a.pdf}{[1]} pdf pages: 3; 4; 5; 6; 7; theorem or section: Assumptions 2.1-2.2; Eq.(11); Definitions 4.1-4.2; Theorems 1-2; theorem: Assumptions 2.1-2.2; Eq.(11); Definitions 4.1-4.2; Theorems 1-2

\needspace{7\baselineskip}\section{References}
[1] Mathieu Dagréou; Thomas Moreau; Samuel Vaiter; Pierre Ablin. A Lower Bound and a Near-Optimal Algorithm for Bilevel Empirical Risk Minimization. 2024. AISTATS 2024, PMLR 238.

\url{https://proceedings.mlr.press/v238/dagreou24a/dagreou24a.pdf}

\end{document}
